% part: intuitionistic-logic
% chapter: introduction
% section: axiomatic-derivations

\documentclass[../../../include/open-logic-section]{subfiles}

\begin{document}

\olfileid{int}{int}{axd}

\olsection{સ્વયંસિદ્ધિમૂલક \usetoken{P}{derivation}}

સ્ફુરણાવાદી વિધાનાત્મક તર્ક માટેની સ્વયંસિદ્ધિમૂલક
!!{derivation}s ખ્યાલની દૃષ્ટિએ સૌથી સરળ અને ઇતિહાસમાં
સૌથી પહેલાં આવેલાં !!{derivation} તંત્રો છે. તેઓ પ્રશિષ્ટ
વિધાનાત્મક તર્કની જેમ જ કાર્ય કરે છે.

\begin{defn}[!!^{derivability}]
જો $\Gamma$ એ $\Lang L$નાં !!{formula}sનો ગણ હોય, તો
$\Gamma$માંથી \emph{!!{derivation}} એ !!{formula}sનો
સાંત ક્રમ $!A_1$, \dots,~$!A_n$ છે, જેમાં દરેક $i \le n$
માટે નીચેનામાંથી કોઈ એક શરત પૂરી થાય છે:
\begin{enumerate}
\item $!A_i \in \Gamma$; અથવા
\item $!A_i$ સ્વયંસિદ્ધિ છે; અથવા
\item $j < i$ અને $k < i$ ધરાવતા કોઈ $!A_j$ અને $!A_k$
  પરથી મોડસ પોનેન્સ વડે $!A_i$ મળે છે, એટલે કે
  $!A_k \ident !A_j \lif !A_i$.
\end{enumerate}
\end{defn}

\begin{defn}[સ્વયંસિદ્ધિઓ]
સ્ફુરણાવાદી વિધાનાત્મક તર્ક માટેની \emph{સ્વયંસિદ્ધિઓ}નો ગણ
$\PAx$ નીચેના સ્વરૂપોનાં બધાં !!{formula}sનો બનેલો છે:
\begin{align}
  & (!A \land !B) \lif !A \ollabel{ax:land1}\\
  & (!A \land !B) \lif !B \ollabel{ax:land2}\\
  & !A \lif (!B \lif (!A \land !B)) \ollabel{ax:land3}\\
  & !A \lif (!A \lor !B) \ollabel{ax:lor1}\\
  & !A \lif (!B \lor !A) \ollabel{ax:lor2}\\
  & (!A \lif !C) \lif ((!B \lif !C) \lif ((!A \lor !B) \lif !C)) \ollabel{ax:lor3}\\
  & !A \lif (!B \lif !A) \ollabel{ax:lif1}\\
  & (!A \lif (!B \lif !C)) \lif ((!A \lif !B) \lif (!A \lif !C)) \ollabel{ax:lif2}\\
  & \lfalse \lif !A \ollabel{ax:lfalse1}
\end{align}
\end{defn}

\begin{defn}[!!^{derivability}]
જો $\Gamma$માંથી અંતે $!A$ આવતું !!a{derivation} હોય, તો
!!^a{formula}~$!A$ એ $\Gamma$માંથી \emph{!!{derivable}}
છે; તેને $\Gamma \Proves !A$ લખીએ છીએ.
\end{defn}

\begin{defn}[પ્રમેયો]
જો ખાલી ગણમાંથી~$!A$નું !!a{derivation} હોય, તો
!!^a{formula}~$!A$ એ \emph{પ્રમેય} છે. $!A$ પ્રમેય
હોય તો $\Proves !A$ અને ન હોય તો $\Proves/ !A$ લખીએ છીએ.
\end{defn}

\begin{prop}
  સ્ફુરણાવાદી તર્કના સ્વયંસિદ્ધિમૂલક !!{derivation} તંત્રમાં
  $\Gamma \Proves !A$ હોય, તો પ્રશિષ્ટ તર્કમાં પણ
  $\Gamma \Proves !A$. ખાસ કરીને, $!A$ સ્ફુરણાવાદી
  પ્રમેય હોય તો પ્રશિષ્ટ પ્રમેય પણ છે.
\end{prop}

\begin{proof}
  દરેક સ્ફુરણાવાદી સ્વયંસિદ્ધિ પ્રશિષ્ટ સ્વયંસિદ્ધિ પણ છે;
  તેથી સ્ફુરણાવાદી તર્કની દરેક !!{derivation} પ્રશિષ્ટ
  તર્કમાં પણ !!a{derivation} છે.
\end{proof}

\end{document}
